% chapters/ch19_parallel.tex
\chapter{Parallel Algorithms for Hierarchical Clustering}
\label{ch:parallel}

\epigraph{%
  Depth is the enemy of speed. Flatten the recursion.%
}{\textit{---Flyxion}}

The top-down recursive algorithm of Chapter~\ref{ch:topdown} has $O(\log n)$ levels of depth, but within each level the split identification tasks for different clusters are entirely independent and can be parallelized. This chapter formalizes the PRAM model and shows that the oracle-augmented partial HC tree construction achieves near-linear total work $W(n) = n^{1+o(1)}$ with only polylogarithmic depth $D(n) = \polylog(n)$. The key observation is that within each cluster, the $O(|S|^2)$ oracle queries are mutually independent and can be distributed across $O(|S|^2)$ processors, collapsing each level's contribution to $O(\log n)$ depth. Over $O(\log n)$ levels the total depth is $O(\log^2 n)$.

%% ── Figure: parallel subtree reconstruction ─────────────────
\begin{figure}[ht]
\centering
\begin{tikzpicture}[
  proc/.style={draw=nodeblue!60, fill=nodeblue!10, rounded corners=3pt,
               minimum width=2.0cm, minimum height=0.7cm, font=\small, align=center},
  arr/.style={->, >=stealth, nodegray!60}
]
% Root split
\node[proc, minimum width=3.5cm] (root) at (0,0) {Split $V$ (parallel queries)};

% Two halves, processed in parallel
\node[proc] (L) at (-2.5,-1.5) {Subtree $A$\\(processor bank 1)};
\node[proc] (R) at (2.5,-1.5)  {Subtree $B$\\(processor bank 2)};
\draw[arr] (root) -- (L);
\draw[arr] (root) -- (R);

% Further splits
\node[proc, minimum width=1.5cm] (LL) at (-3.8,-3.0) {$A_1$};
\node[proc, minimum width=1.5cm] (LR) at (-1.2,-3.0) {$A_2$};
\node[proc, minimum width=1.5cm] (RL) at (1.2,-3.0) {$B_1$};
\node[proc, minimum width=1.5cm] (RR) at (3.8,-3.0) {$B_2$};
\draw[arr] (L)--(LL); \draw[arr] (L)--(LR);
\draw[arr] (R)--(RL); \draw[arr] (R)--(RR);

% Timing annotation
\draw[decorate, decoration={brace, amplitude=5pt, mirror}, nodegreen!60]
  (-4.8,-3.5) -- (4.8,-3.5);
\node[font=\footnotesize, nodegreen!70, below] at (0,-3.9)
  {all four splits at same depth completed in $O(\log n)$ time};

% Work annotation
\node[font=\footnotesize, nodered!70, right] at (5.0,0)
  {$W=n^{1+o(1)}$};
\node[font=\footnotesize, nodered!70, right] at (5.0,-1.5)
  {$D=\polylog(n)$};
\end{tikzpicture}
\caption{Parallel reconstruction: different subtrees are processed by independent processor banks simultaneously. Within each subtree, oracle queries are also parallelized. The result is a polylogarithmic-depth algorithm with near-linear total work.}
\label{fig:parallel}
\end{figure}

\section{The PRAM Model}

We measure \emph{work} $W(n)$ (total operations) and \emph{depth} $D(n)$ (parallel time, i.e., the longest sequential dependency chain).

\section{Main Result}

\begin{theorem}[Parallel MW approximation]
\label{thm:parallel}
There exists a PRAM algorithm for the Moseley--Wang objective with
\[
W(n) = n^{1+o(1)}, \qquad D(n) = \operatorname{polylog}(n),
\]
achieving a $(1-o(1))$-approximation.
\end{theorem}

\section{Exercises}

\begin{enumerate}
  \item Show the top-down split identification is embarrassingly parallelizable and compute the depth per level.
  \item Compare the work of the parallel algorithm with the sequential algorithm.
  \item Describe how the parallel algorithm handles load balancing for unequal cluster sizes.
\end{enumerate}
